\chapter{Revisión bibliográfica: construcción de la base documental}\label{cap:revision}

Este capítulo desarrolla la \textbf{Fase 1 de la metodología}: describe y justifica el procedimiento completo de la revisión (búsqueda, cribado, extracción y síntesis) y presenta sus resultados, de modo que quede trazable cómo se construyó la base documental que sustenta la propuesta didáctica del Capítulo 3.

\section{Objetivo y preguntas de la revisión}\label{sec:rev-objetivo}

La revisión persigue fundamentar el diseño de una secuencia didáctica basada en el modelo 5E para \textbf{la determinación de la edad y el tamaño del universo} en 2º de Bachillerato (con la expansión y la ley de Hubble-Lemaître como vía de medida, y la expansión acelerada y la energía oscura como matiz). Para ello se formulan tres preguntas de revisión:

\begin{itemize}
\item \textbf{PR1.} ¿Qué contenidos de cosmología moderna se han trabajado con mayor frecuencia en bachillerato o niveles equivalentes?
\item \textbf{PR2.} ¿Qué estrategias didácticas o de divulgación se han empleado para enseñarlos o comunicarlos?
\item \textbf{PR3.} ¿Qué dificultades conceptuales e ideas erróneas se han documentado?
\end{itemize}

De forma transversal, cada fuente se analiza además por su \textbf{utilidad para el diseño de la secuencia 5E} (elementos transferibles y fase del modelo en la que encajarían).

\section{Metodología de la revisión}\label{sec:rev-metodologia}

Se siguió un procedimiento de revisión exploratoria sistematizada (tipo \textit{scoping review}), apoyado en la herramienta \textbf{Elicit Pro}, y estructurado en cuatro etapas: búsqueda, cribado (\textit{screening}), extracción de datos y síntesis. La selección de fuentes se documenta siguiendo la lógica de \textbf{PRISMA-ScR} (extensión de PRISMA para revisiones exploratorias), reportando en el texto los recuentos de cada etapa: identificación, cribado con motivos de exclusión e inclusión.

\subsection{Estrategia de búsqueda}

Se realizó una búsqueda semántica sobre el índice de Elicit ($>$138 millones de artículos, que integra Semantic Scholar y OpenAlex) con la consulta: \textbf{``What educational strategies have been used to teach or communicate dark energy, cosmic expansion, and modern cosmology to high school students in scientifically accurate and accessible ways?''}. La búsqueda devolvió \textbf{1.000 resultados}, que constituyeron el conjunto inicial para cribado. El rango temporal del corpus resultante abarca \textbf{1991--2026}.

\subsection{Criterios de cribado (\textit{screening})}

El cribado por resúmenes aplicó criterios relativos a: \textbf{(i)} población destinataria (alumnado, profesorado, currículos o materiales para \textasciitilde14--18 años); \textbf{(ii)} pertinencia del contenido de cosmología moderna (Big Bang, expansión, redshift, ley de Hubble-Lemaître, CMB, materia oscura, energía oscura, aceleración); \textbf{(iii)} enfoque educativo o divulgativo; \textbf{(iv)} calidad de la fuente (revisión por pares, actas, libros académicos o repositorios reputados); y \textbf{(v)} nivel educativo apropiado (secundaria/bachillerato o divulgación adaptada, no exclusivamente primaria o universidad). Las fuentes que incumplían algún criterio estricto se excluyeron automáticamente; el resto se valoró de forma holística.

De los 1.000 registros identificados se excluyeron \textbf{910} en el cribado, con la siguiente distribución por motivo principal: pertinencia del contenido cosmológico = 487; nivel/edad no apropiados = 330; precisión conceptual insuficiente = 35; ausencia de dificultades/concepciones = 23; contenido no sustantivo = 14; y motivos menores = 21. Tras el cribado quedaron \textbf{90 fuentes incluidas} (1.000 $-$ 910 = 90).

\subsection{Extracción de datos}

De cada fuente incluida se extrajo información en \textbf{diez dimensiones}, que después alimentan la tabla de análisis y la síntesis: (1) estrategia educativa; (2) contenido de cosmología; (3) contexto de implementación; (4) precisión científica; (5) accesibilidad conceptual; (6) resultados de aprendizaje; (7) motivación/\textit{engagement}; (8) dificultades identificadas; (9) concepciones alternativas; y (10) elementos transferibles al diseño de la secuencia. Para las fuentes disponibles solo como resumen, la extracción se limitó a la información del \textit{abstract} (véase la sección de limitaciones).

\subsection{Depuración, síntesis y categorización}

Sobre las 90 fuentes incluidas se aplicaron tres operaciones:

\begin{enumerate}
\item \textbf{Fusión de duplicados.} La revisión identifica \textbf{ocho pares} que reportan el mismo estudio o recurso (Etkina et al., 1999; Pimbblet \& Newman, 2003; Çoban \& Sengoren, 2011; Navone et al., 2011; Bartlett, 2018; Tuveri et al., 2022; Cardona-Rodríguez et al., 2016; y Pereira, 2019). Cada par se consolida en una sola entrada, resultando \textbf{82 fuentes únicas}.
\item \textbf{Exclusión por baja pertinencia.} Se descartaron \textbf{cuatro} fuentes (solo resumen) por contenido cosmológico marginal o aporte insuficiente (Adam, 2016; Jacobs, 2013; Arcadias et al., 2025; y Ritz \& Chow, 2018). Se \textbf{conservaron} Etkina et al. (1999) ---pese a su contenido más estelar que cosmológico--- y Novotný \& Svobodová (2017) ---pese a su nivel universitario---, dejando constancia de ambas cautelas.
\item \textbf{Categorización temática.} Las \textbf{78 fuentes} resultantes se agruparon en \textbf{diez categorías de estrategia/contexto}, que ordenan tanto la síntesis como la tabla de análisis.
\end{enumerate}

El conjunto del proceso de selección se resume en la Figura~\ref{fig:prisma}.

\begin{figure}[htbp]
\centering
\begin{tikzpicture}[
  font=\footnotesize,
  node distance=6mm and 9mm,
  caja/.style={rectangle, rounded corners, draw=naranja, line width=0.7pt, fill=naranja!10, text width=6.0cm, inner sep=2mm, align=center},
  excl/.style={rectangle, rounded corners, draw=black!30, dashed, fill=black!5, text width=5.2cm, inner sep=1.6mm, align=left, font=\scriptsize},
  flecha/.style={-{Stealth[length=2.2mm]}, draw=naranja, line width=0.9pt},
  fexcl/.style={-{Stealth[length=1.8mm]}, draw=black!35, dashed, line width=0.6pt}
]
\node[caja] (id) {\textbf{Identificación.} Registros identificados con Elicit Pro: \textbf{$n=1000$}};
\node[caja, below=of id] (inc) {\textbf{Cribado por título y resumen.} Fuentes incluidas: \textbf{$n=90$}};
\node[caja, below=of inc] (uni) {Fuentes \textbf{únicas} tras fusionar duplicados: \textbf{$n=82$}};
\node[caja, below=of uni] (ana) {\textbf{Fuentes analizadas} (diez categorías): \textbf{$n=78$}};
\node[excl, right=of inc] (e1) {\textbf{Excluidos en el cribado: $n=910$.} Pertinencia del contenido (487); nivel/edad (330); precisión conceptual (35); sin dificultades/concepciones (23); no sustantivo (14); otros (21).};
\node[excl, right=of uni] (e2) {$-8$ pares \textbf{duplicados} fusionados.};
\node[excl, right=of ana] (e3) {$-4$ fuentes por \textbf{baja pertinencia} (solo resumen).};
\draw[flecha] (id) -- (inc);
\draw[flecha] (inc) -- (uni);
\draw[flecha] (uni) -- (ana);
\draw[fexcl] (inc) -- (e1);
\draw[fexcl] (uni) -- (e2);
\draw[fexcl] (ana) -- (e3);
\end{tikzpicture}
\caption{Flujo de selección de fuentes (PRISMA-ScR): de los 1000 registros identificados a las 78 fuentes analizadas.}\label{fig:prisma}
\end{figure}

\noindent\textit{Convención de la tabla.} Las fuentes disponibles solo como resumen se marcan con \dag{}; en ellas las columnas de dificultades y elementos 5E son más tentativas. La columna ``Fase 5E'' indica la fase del modelo en la que el elemento resultaría más útil; para las fuentes de contexto o metateóricas se indica ``Base de diseño'' o ``Contexto/diseño''. El término ``pre-post'' designa la medición de la comprensión antes y después de la intervención. Los índices [n] remiten a la numeración del listado de fuentes identificadas con Elicit Pro, incluido como anexo.

\section{Resultados}\label{sec:rev-resultados}

\subsection{Características del corpus}

Las 78 fuentes analizadas ---de las 82 únicas, tras descartar cuatro por baja pertinencia--- abarcan 1991--2026 e incluyen intervenciones de aula, análisis curriculares, programas de extensión, revisiones conceptuales y descripciones de recursos. Proceden de contextos diversos (Australia, Brasil, Argentina, Colombia, Alemania, Italia, España, Suecia, Estados Unidos, Indonesia, Taiwán, Ucrania y Reino Unido), con predominio de Brasil, Alemania, Australia y Reino Unido. Aproximadamente \textbf{la mitad} dispone de texto completo; la otra mitad solo de resumen. Un patrón recurrente es su origen en talleres de profesorado, tesis, actas de congresos y revistas de práctica docente, más que en estudios experimentales controlados, lo que condiciona la fuerza de la evidencia sobre eficacia.

\subsection{PR1 --- Temas de cosmología moderna más frecuentes}\label{sec:pr1}

Ordenados de mayor a menor presencia en el corpus:

\begin{enumerate}
\item \textbf{Expansión del universo y ley de Hubble-Lemaître} --- columna vertebral de la revisión; el tema más trabajado y transferible.
\item \textbf{Corrimiento al rojo (redshift)} --- casi siempre ligado a la expansión; abordado tanto conceptual (Doppler vs. cosmológico) como cuantitativamente.
\item \textbf{Big Bang y evolución del universo} --- muy frecuente, sobre todo en enfoques históricos y narrativos.
\item \textbf{Materia oscura} --- muy presente en divulgación y \textit{outreach} (masterclasses, lentes gravitacionales, curvas de rotación).
\item \textbf{Radiación cósmica de fondo (CMB)} --- como evidencia observacional más que como actividad autónoma.
\item \textbf{Galaxias y estructura a gran escala} --- soporte observacional.
\item \textbf{Energía oscura y aceleración cósmica} --- \textbf{uno de los temas menos tratados en profundidad}; suele aparecer como punto final de un relato, y es el más vulnerable a presentarse como ``sustancia conocida''.
\item \textbf{Edad del universo} --- subtema derivado de forma natural de $H_0$ ($t \approx 1/H_0$), apenas tratado explícitamente \parencite{carmesincarmesin2014,linton2012}: \textbf{oportunidad de originalidad} para el TFM.
\item \textbf{Relatividad general y espacio-tiempo} --- reservado a cursos avanzados o alumnado de altas capacidades.
\end{enumerate}

\subsection{PR2 --- Estrategias didácticas y de divulgación (diez categorías)}\label{sec:pr2}

\begin{table}[ht]
\centering
\begin{tabular}{|c|l|c|} \hline
\rowcolor{naranja}
{\color{white}\textbf{\#}} & {\color{white}\textbf{Categoría}} & {\color{white}\textbf{N.º fuentes}} \\ \hline
1 & Analogías y modelos físicos & 11 \\ \hline
2 & Análisis de datos reales / actividades con datos & 7 \\ \hline
3 & Simulaciones y visualizaciones digitales & 2 \\ \hline
4 & Modelado matemático formal & 12 \\ \hline
5 & Secuencias didácticas e indagación en el aula (ABP/TBL/aula invertida) & 6 \\ \hline
6 & Programas de extensión, inmersión y \textit{masterclasses} & 12 \\ \hline
7 & Enfoques históricos y de naturaleza de la ciencia & 4 \\ \hline
8 & Recursos lúdicos, artísticos y de divulgación & 8 \\ \hline
9 & Instrumentos de evaluación y diagnóstico de concepciones & 7 \\ \hline
10 & Currículo, formación del profesorado y barreras contextuales & 9 \\ \hline
\end{tabular}
\end{table}

Las \textbf{analogías físicas} (globo, banda elástica, ondas de agua, lentes) son la estrategia más común, pero su eficacia depende críticamente de que \textbf{se expliciten sus límites}; el \textbf{análisis de datos reales} aparece como la vía más honesta científicamente para llegar a la energía oscura; y los \textbf{programas de extensión} concentran buena parte del trabajo sobre materia oscura. Las categorías 9 y 10 no son ``estrategias'' en sentido estricto, pero son imprescindibles para el diseño: aportan los \textbf{instrumentos de evaluación} y documentan las \textbf{barreras} (formación del profesorado, tiempo, currículo).

\subsection{PR3 --- Dificultades conceptuales e ideas erróneas}\label{sec:pr3}

El corpus converge en un conjunto consistente entre países y edades:

\begin{enumerate}
\item \textbf{Big Bang como explosión en un punto} del espacio vacío (en vez de expansión del propio espacio).
\item \textbf{El universo tiene centro y borde.}
\item \textbf{Los sistemas ligados} (galaxias, sistemas solares, átomos) \textbf{se expanden} con el espacio.
\item \textbf{Todo corrimiento al rojo es efecto Doppler}, en lugar del estiramiento de la longitud de onda.
\item \textbf{Confusión materia oscura / energía oscura} (y de la materia oscura con antimateria y agujeros negros).
\item \textbf{Energía oscura como fuerza simple que ``empuja''} a las galaxias.
\item \textbf{``Expansión'' como creación de materia nueva} o como ``nuevos descubrimientos''.
\item \textbf{La luz de objetos lejanos no se ve afectada} por la expansión.
\item \textbf{No conservación de la energía del fotón} que se corre al rojo.
\item \textbf{Conflicto lenguaje cotidiano / técnico} en torno a la palabra ``expansión''.
\end{enumerate}

\noindent De estas diez dificultades, las \textbf{seis primeras} constituyen el \textbf{núcleo} que la secuencia confronta de forma explícita ---las «seis ideas erróneas nucleares» a las que remiten los Capítulos 3 y 4---; las cuatro restantes se atienden de forma secundaria.

A ello se suman \textbf{barreras sistémicas}: déficit de formación del profesorado, saturación curricular y escaso tiempo lectivo. La evidencia \parencite{aretz2017} advierte de que estas concepciones \textbf{conviven con altos niveles de conocimiento factual}, por lo que su corrección exige actividades de cambio conceptual dirigidas, no mera transmisión de información.

\subsection{Tabla de análisis por fuente}\label{sec:tabla}

La Tabla~\ref{tab:analisis} recoge las 78 fuentes agrupadas en las diez categorías anteriores. Los índices [n] remiten al listado de fuentes (anexo).

\begin{landscape}
\footnotesize
\setlength{\tabcolsep}{2pt}
\input{capitulos/tabla_analisis}
\end{landscape}

\section{Discusión}\label{sec:discusion}

\subsection{¿Qué temas son más viables para el bachillerato?}

El bloque \textbf{expansión -- Hubble-Lemaître -- redshift} es el más viable: admite tratamientos desde lo cualitativo (analogías) hasta lo cuantitativo con matemática elemental (ajuste de una recta $v$--$d$ y cálculo de $H_0$). De él deriva la \textbf{edad del universo} ($t \approx 1/H_0$), objetivo concreto, medible y poco explotado en la literatura. El \textbf{Big Bang} y la \textbf{CMB} son viables sobre todo como parte de un relato histórico-observacional. La \textbf{energía oscura y la aceleración}, aun siendo el tema ``estrella'', son los menos tratados en profundidad y los más expuestos a simplificaciones, por lo que conviene situarlos como \textbf{matiz o extensión hacia el final} de la secuencia, no como núcleo ni como entrada. La relatividad general formal y las distancias cosmológicas quedan para alumnado avanzado.

\subsection{¿Cómo introducir la energía oscura y la expansión acelerada como matiz final?}\label{sec:disc-estrategias}

La vía más honesta es la \textbf{secuencia con datos reales} de \textcite{hyde2021}: reproducir el diagrama de Hubble, prever la desaceleración por gravedad y, con los datos de supernovas, observar una \textbf{anomalía} (la aceleración) que solo se explica introduciendo la energía oscura ---presentada así como el concepto que mejor da cuenta de lo observado, no como sustancia---. Se apoya en simuladores que la hacen manipulable \parencite{raymundo2022} y en la \textbf{paradoja newtoniana} de \textcite{pereira2019}, que la ancla en la conservación de la energía evitando la imagen de ``fuerza que empuja''. Como andamiaje son útiles las \textbf{analogías con límites explícitos} \parencite{xu2018,suarez2025} y la demostración de \textcite{marshall2022} para separar el corrimiento al rojo Doppler y el cosmológico. La lección transversal: la analogía no es buena o mala en sí, sino según \textbf{se expliciten sus límites}.

\subsection{¿Qué dificultades deben abordarse desde el diseño?}

La secuencia debe atacar explícitamente las seis ideas erróneas nucleares, anticipar el conflicto \textbf{lenguaje cotidiano/técnico} en ``expansión'' \parencite{salimpour2023} y asumir que el cambio conceptual requiere \textbf{confrontar la preconcepción}, no solo informar \parencite{aretz2017}.

\subsection{¿Qué elementos encajan en cada fase 5E?}\label{sec:disc-5e}

\begin{itemize}
\item \textbf{Engage:} paradoja de Olbers \parencite{santos2014}; narrativa histórica \parencite{lochner2009,nunes2025}; diagnóstico de concepciones \parencite{colantonio2019}.
\item \textbf{Explore:} medición $v$--$d$ con banda elástica y Tracker \parencite{suarez2025} o con datos reales/SDSS \parencite{hyde2021,cardona2016}; obtención de $H_0$.
\item \textbf{Explain:} de $H_0$ a la edad $t \approx 1/H_0$; analogía del globo con límites \parencite{xu2018}; Doppler vs. cosmológico con coches RC \parencite{marshall2022}; marco de \textcite{possel2020}.
\item \textbf{Elaborate:} supernovas $\rightarrow$ aceleración $\rightarrow$ energía oscura \parencite{hyde2021,raymundo2022}; paradoja newtoniana \parencite{pereira2019}; tamaño y edad del universo observable \parencite{linton2012}.
\item \textbf{Evaluate:} ítems adaptados de CosmoCI \parencite{salimpour2022} y del \textit{construct map} de \textcite{aretz2017}; \textcite{gasparini2025} como referente para el futuro diseño pre-post.
\end{itemize}

\section{Anclaje curricular en el contexto de implementación}\label{sec:anclaje}

Para valorar la viabilidad real de la propuesta se cotejan los contenidos de la revisión con los currículos de Física que imparte el autor: \textbf{Pearson Edexcel International Advanced Level (IAL) Physics} (años 1 y 2, equivalentes a 1º y 2º de Bachillerato) y la \textbf{PCE de Física de la UNED} (2º de Bachillerato, RD 243/2022, LOMLOE).

\subsection{Correspondencia de contenidos}

\begin{table}[ht]
\centering
\footnotesize
\begin{tabularx}{\textwidth}{|X|c|c|c|} \hline
\rowcolor{naranja}
{\color{white}\textbf{Contenido de la revisión}} & {\color{white}\textbf{Edexcel Año 1}} & {\color{white}\textbf{Edexcel Año 2 (T11)}} & {\color{white}\textbf{PCE UNED}} \\ \hline
Campo gravitatorio, potencial, órbitas & Parcial & \checkmark\ (11A) & \checkmark\ (Bloque II) \\ \hline
Ondas y efecto Doppler & Base & \checkmark & \checkmark\ (Ondas) \\ \hline
Corrimiento al rojo ($z$, $v = H_0 d$) & --- & \checkmark & Parcial \\ \hline
Ley de Hubble-Lemaître / $H_0$ & --- & \checkmark & --- \\ \hline
\textbf{Edad del universo} & --- & \checkmark & Implícito \\ \hline
Expansión y \textbf{destino del universo} & --- & \checkmark & --- \\ \hline
\textbf{Energía oscura} & --- & \checkmark & --- \\ \hline
\textbf{Materia oscura} & --- & \checkmark & --- \\ \hline
Big Bang y CMB & --- & \checkmark & \checkmark\ (Bloque VI) \\ \hline
Clasificación estelar, H-R, distancias & --- & \checkmark & --- \\ \hline
Relatividad especial & --- & --- & \checkmark\ (Bloque VI) \\ \hline
\end{tabularx}
\end{table}

\subsection{Edexcel IAL Physics}

El \textbf{Año 1} no contiene cosmología, pero aporta prerrequisitos clave (concepto de onda, conservación de la energía, energías cinética y potencial gravitatoria). El \textbf{Año 2, Topic 11 ``Astrophysics and Cosmology''}, encaja de forma casi literal con el objeto del TFM: la especificación (unidad 5.6) exige campo gravitatorio y órbitas, radiación de cuerpo negro (Stefan-Boltzmann, Wien), distancias por paralaje y candelas estándar, diagrama de Hertzsprung-Russell, efecto Doppler y corrimiento al rojo ($z = \Delta\lambda/\lambda \approx v/c$ y $v = H_0 d$), y explícitamente la controversia sobre la edad y el destino del universo asociada al valor de la constante de Hubble y a la posible existencia de materia oscura. El libro de texto desarrolla las secciones ``The Age of the Universe'', ``The Fate of the Universe'' y un \textit{Thinking Bigger} sobre energía oscura. Así, la secuencia propuesta ---expansión $\rightarrow$ ley de Hubble-Lemaître $\rightarrow$ edad del universo--- \textbf{coincide con los estándares evaluables de la Unit~5.6} (subtema 11B ``Space'' del libro), con un enfoque coherente con la revisión \parencite{hyde2021,pimbblet2003}.

\subsection{PCE UNED (LOMLOE)}

La cosmología aparece más repartida y modesta: el \textbf{Bloque II (Interacción gravitatoria)} cubre campo, potencial y movimiento orbital; el \textbf{Bloque VI (Física del siglo XX)} incluye ``Historia y composición del Universo'' y exige explicar la teoría del Big Bang y sus evidencias (radiación de fondo y efecto Doppler relativista). El Big Bang, la CMB y el Doppler relativista son, pues, contenidos exigibles, pero \textbf{la edad del universo, la ley de Hubble-Lemaître cuantitativa y la energía oscura no figuran de forma explícita}: en este marco la secuencia funcionaría como \textbf{ampliación} del Bloque VI, apoyada en la gravitación del Bloque II.

\subsection{Conclusión sobre el nivel y el vehículo de implementación}

El cotejo refuerza tres decisiones: \textbf{(i)} el nivel idóneo es \textbf{2º de Bachillerato / Año 2 de Edexcel}, donde los contenidos son currículo evaluable, no un añadido; \textbf{(ii)} el \textbf{vehículo de implementación} más directo es la \textit{Unit~5.6} de Edexcel (subtema 11B ``Space'' del libro), con la PCE UNED como marco alternativo en clave de ampliación; y \textbf{(iii)} presentar la secuencia 5E como una forma de \textbf{enseñar contenido ya prescrito} neutraliza la objeción de ``saturación curricular'' que la literatura identifica como principal barrera \parencite{skolimoski2014,lin2011}.

\section{Implicaciones para el diseño de la secuencia 5E}\label{sec:implicaciones}

De la revisión se derivan cinco decisiones de diseño: \textbf{(i)} adoptar una \textbf{secuenciación observacional} (Hubble $\rightarrow$ deceleración $\rightarrow$ supernovas $\rightarrow$ energía oscura), coherente con \textcite{pimbblet2003} e \textcite{hyde2021}; \textbf{(ii)} usar analogías \textbf{solo con sus límites explicitados}; \textbf{(iii)} situar la determinación de la \textbf{edad} del universo como resultado tangible de la medida (fases Explore y Explain) y la del \textbf{tamaño} observable como su extensión directa (Elaborate), con la \textbf{energía oscura} como matiz que refina y acota la estimación (también en Elaborate); \textbf{(iv)} abordar de forma explícita las seis ideas erróneas nucleares; y \textbf{(v)} diseñar la evaluación con instrumentos validados (CosmoCI, \textit{construct map}), reservando su aplicación pre-post para la implementación en el curso 2026-2027.

\section{Limitaciones de la revisión}\label{sec:limitaciones}

\begin{itemize}
\item \textbf{Cobertura de una sola fuente de búsqueda} (Elicit Pro): puede faltar literatura no indexada o en repositorios locales.
\item \textbf{Aproximadamente la mitad del corpus solo dispone de resumen} (fuentes marcadas con \dag{}), lo que limita la profundidad de la extracción en esas entradas.
\item \textbf{Escasez de estudios con diseño pre-post e instrumentos validados}: solo \textcite{gasparini2025} y \textcite{dua2020} permiten estimar tamaños del efecto, lo que impide identificar una única estrategia óptima.
\item \textbf{Sesgo de idioma y contexto}: predominio de fuentes de Brasil, Alemania, Australia y Reino Unido; el encaje con los currículos del contexto de implementación se aborda en la sección de anclaje curricular.
\end{itemize}

La revisión de 78 fuentes (1991--2026, de las 82 únicas del corpus) confirma que la \textbf{expansión del universo} es el contenido más trabajado y viable, que las \textbf{analogías con límites explícitos} y el \textbf{análisis de datos reales} son las estrategias más fértiles ---esta última, la vía más honesta hacia la energía oscura---, y que existe un \textbf{núcleo estable de ideas erróneas} que la secuencia debe atacar de forma dirigida. Estos resultados fundamentan el diseño de la secuencia 5E del Capítulo 3, con la determinación de la \textbf{edad y el tamaño del universo} como objetivo central y la \textbf{energía oscura} como matiz.